% == == == == == == == == == == == == == == == == == == == == == == == == == ==
    == == == == == == == == == == == ==
    = % Section 14 : Lessons Learned % == == == == == == == == == == == == == ==
      == == == == == == == == == == == == == == == == == == == == == == == ==
    =

\section {
  Lessons Learned
}
\sectionauthors{CM, BS, JH, MN, SC, JB} %
    Reviewed by Jeremie

        This section documents key insights gained during the project lifecycle,
    demonstrating continuous improvement and knowledge capture
        .These lessons have been applied iteratively throughout the project
            phases.

\subsection {
  Research Phase Lessons
}

\subsubsection { LL - 1 : Iterative Compute Platform Selection }

\textbf{Situation}: During compute platform selection, the team evaluated multiple options for the main processing unit. The initial candidate was the Texas Instruments SK-TDA4VM, selected for its integrated vision accelerators and automotive-grade reliability.

\textbf{
  Discovery \#1
}: Power analysis revealed the SK-TDA4VM consumed 15--20W under typical SLAM workloads and 100W at max load, exceeding our 10W power budget for the compute subsystem. This would have required a larger battery and increased robot weight.

\textbf{
  Pivot \#1
}: Research into low-power FPGA alternatives led to the Pynq-Z2 board, inspired by work on unified SLAM accelerators for low-end FPGAs \parencite{sugiura_fpga_slam}. The Pynq-Z2 offered hardware acceleration within our power envelope.

\textbf{
  Discovery \#2
}: Compatibility testing revealed that ROS 2 Humble had limited support for the Pynq-Z2's ARM architecture and required significant manual configuration. The FPGA development toolchain also added complexity.

\textbf{
  Pivot \#2
}: The team selected Raspberry Pi 5, which provides native Ubuntu 22.04 and ROS 2 Humble support, extensive community documentation, proven reliability in robotics applications, and acceptable power consumption (5--8W under load).

\textbf{Result}: Zero schedule impact. Both pivots occurred during Research before any hardware procurement. Software development proceeded without compatibility issues.

\textbf{Lesson}: \textit{
  Evaluate multiple candidates against all constraints(
      power, software compatibility, development complexity)
      before procurement.Early pivots cost hours;
  late pivots cost weeks.
}

\subsubsection { LL - 2 : Sponsor Relationship Investment }

\textbf{Situation}
    : The project budget constrained sensor options,
      with professional - grade LiDARs costing \$2,
      000 --\$5,
      000.

\textbf{Action}
    : Proactive engagement with SICK Inc.through multiple meetings with
          Dr.Porter clarified project scope and demonstrated educational value
              .The team presented the
                  project 's alignment with SICK' s educational outreach goals
              .

\textbf{Result} : SICK Inc.donated a TiM561 LiDAR valued at \$3,
      500 +
    ,
      providing professional - grade sensors at educational budget levels.

\textbf{Lesson} : \textit {
  Sponsor relationships provide resources beyond funding
      : expertise,
        equipment,
        and scope guidance.Invest time in demonstrating mutual value.
}

\subsection { Design Phase Lessons }

\subsubsection { LL - 3 : Multi - Board PCB Strategy }

\textbf{Situation}: The original electrical design called for a single large PCB containing all motor drivers, power distribution, and sensor interfaces.

\textbf{Risk Analysis}: Risk assessment revealed that a single board failure would be catastrophic;
the entire system would be non - functional,
    and replacement would require complete board redesign and fabrication

\textbf{Action} : Redesigned to six smaller PCBs.Each board uses 2 - layer or
        4 - layer and can be replaced independently.

\textbf{Result} : If one board fails,
    it's cheaper to replace one board than the entire system.

\textbf{Lesson} : \textit {Design for failure isolation. Modular hardware architecture reduces both cost and catastrophic failure risk.
}

\subsubsection {LL-4: Interim Hardware for Parallel Development
}

\textbf{Situation}
    : The donated SICK TiM561 LiDAR
        requires a
      specific power adapter board.Initial testing revealed we did not possess
          the correct adapter,
      and sourcing would take 2 --3 weeks.

\textbf{Problem} : Without a functional LiDAR,
      software
      development(ROS 2 drivers, SLAM integration, visualization) would stall.

\textbf{Action}: Ordered an RPLiDAR C1 (\$55) as an interim development sensor. Both sensors use similar ROS 2 driver interfaces, allowing software development to proceed.

\textbf{Result}: Software team continued SLAM development and ROS 2 integration while the SICK adapter was sourced. When the correct adapter arrived, the software was ready for immediate integration.

\textbf{Lesson}: \textit{
  Identify blockers early and find workarounds to maintain parallel
      progress.A \$55 interim sensor prevented weeks of idle time.
}

\subsection { Implementation Insights }

\subsubsection { LL - 5 : Backup Hardware Strategy Validated }

\textbf{Incident}
    : During initial testing in November 2025,
      jumper wires connected to the SICK TiM561 LiDAR melted due to high inrush
              current at startup.The
                  LiDAR draws significant current during power -
          on,
      which exceeded the current capacity of the temporary jumper wire
              connections
                  .

\textbf{Impact} : Initial testing was temporarily halted while the
                       team investigated the root cause.Testing with a current -
          limited bench power supply confirmed the LiDAR itself
                  was undamaged and
                      operated normally under controlled conditions.

\textbf{Mitigation} : Because the RPLiDAR C1 had been ordered as a backup /
              development sensor(per LL - 4),
      software development continued uninterrupted using the backup sensor
              .The team was able to proceed with SLAM development and ROS
          2 integration while addressing the power delivery issue
              .

\textbf{Resolution}
    : The SICK TiM561 remains the primary sensor with proper power
      management.The RPLiDAR C1 continues to serve as the backup
          / development sensor,
      providing redundancy and enabling parallel testing.

\textbf{Lesson} : \textit {Always have backup hardware for critical components, and test power delivery carefully before connecting sensitive equipment. Inrush current management is critical for industrial sensors.
}

\subsection{Applied Improvements Summary}

These lessons were applied iteratively throughout the project :


\begin{table}[H]
\centering
\caption {
  Lessons Applied During Project
}
\label {
tab:
  lessons_applied
}
\begin{tabular} {
  @{} llll @{}
}
\toprule
\textbf{Lesson} & \textbf{Phase Applied} & \textbf{Action Taken} & \textbf {
  Outcome
}
\\
\midrule LL - 1 &
    Research &Pynq - Z2 $\rightarrow$ RPi 5 pivot &Zero schedule impact \\ 
LL - 2 &
    Research &SICK sponsorship engagement & \$3,
    500 + equipment donated \ LL - 3 &
        Design &Single PCB $\rightarrow$ 6 boards &Cost reduction \ LL - 4 &
        Design &Ordered RPLiDAR C1 interim &Parallel development \ LL - 5 &
        Build(early) & Used backup during power issue &No schedule slip \\
\bottomrule
\end {
  tabular
}
\end { table }

\subsection {Recommendations for Future Projects
}

Based on these lessons, we recommend :

\begin {
  enumerate
}
\item \textbf{Validate Early} : Test software /
                                hardware compatibility during Research,
    not Build
    \item \textbf {Design for Failure
}: Modular architectures enable graceful degradation
    \item \textbf{Budget for Backups
}: Include backup components for critical-path items
    \item \textbf{Engage Sponsors}: Industry partnerships provide more than funding
\end{
  enumerate
}
